\chapter{Exercício 9.2}

Todo o código usado nesse exercício encontra-se no Anexo \ref{cap:ex9-2-codigo}.

Considerando o problema do exercício 8.2 em que temos um domínio 2D quadrado e 
queremos simular o modo $TM_z$ com o FDTD, usamos os seguintes parâmetros de discretização:

\begin{itemize}
	\item $N_x = 200 \quad , \quad N_y = 200$
	\item $\Delta y = \Delta x = 1 \un{mm}$
	\item $\Delta t = \frac{0.9}{c \sqrt{(\Delta x)^{2} + (\Delta y)^{2}}} = 2.12 \un{ps} $ (condição CFL 2D)
	\item $N_t = 800$
\end{itemize}

A fonte aplicada no centro do domínio é uma função gaussiana dada por 

\[ f(t) = e^{-\frac{1}{2}\left(\frac{n - 40}{12}\right)^2} \]

\noindent que tem um pico em $t = 84 \un{ps}$ e uma largura de aproximadamente $25 \un{ps}$.

Os parâmetros da CPML adotados foram os seguintes:

\begin{itemize}
	\item $N_{pml} = 30 \quad , \quad m = 3$
	\item $R_{err} = 10^{-6}$
	\item $\eta = \sqrt{\frac{\mu}{\epsilon}}$
	\item $\sigma_{max} = -\frac{(m+1)\ln(R_{\mathrm{err}})}{2\eta N_{pml}} = 2.44$
\end{itemize}

A propagação obtida está exibida na Figura \ref{fig:ex9-2-propagacao}, 
onde podemos observar que o campo se propaga para fora do domínio sem reflexões significativas, 
indicando que os PMLs estão funcionando corretamente. Com uma ponta de prova
próximo ao centro do grid, verificamos uma leve reflexão do campo,
como mostra a Figura \ref{fig:ex9-2-reflexao}.

\begin{figure}[htbp]
	\caption{\label{fig:ex9-2-propagacao}Propagação de $E_z$ ao longo do grid, sem reflexão nas bordas.}
	\begin{center}
        \includegraphics[width=\textwidth,trim=1 1 1 1,clip]{ex9-2-propagacao.png}
	\end{center}
\end{figure}

\begin{figure}[htbp]
	\caption{\label{fig:ex9-2-reflexao}Reflexão de $E_z$ no centro do grid em função do tempo.}
	\begin{center}
        \includegraphics[width=\textwidth,trim=1 1 1 1,clip]{ex9-2-reflexao.png}
	\end{center}
\end{figure}

O coeficiente de reflexão pode ser calculado a partir dos picos verificados na Figura \ref{fig:ex9-2-reflexao}, 
obtendo: 

\[ R = \frac{0.00007}{0.01} = 0.007 = 0.7 \% \]

\noindent o qual, quando comparado com a reflexão obtida usando o Mur de $R = 5 \%$ no exercício 8.2, 
representa uma redução quase 10 vezes melhor usando a PML.

A Figura \ref{fig:ex9-2-propagacao-x} mostra a evolução do campo $E_z$ ao longo do tempo, em um corte horizontal no eixo x
passando pelo centro da cavidade.
 As linhas verticais vermelhas indicam a região onde o campo é absorvido pelos PMLs, e podemos 
 observar que o campo decai para zero em todos os instantes, evidenciando a absorção da PML.

\begin{figure}[htbp]
	\caption{\label{fig:ex9-2-propagacao-x}Decaimento do campo $E_z$ ao longo do eixo $x$ ao entrar nas regiões de PML.}
	\begin{center}
        \includegraphics[width=0.8\textwidth,trim=1 1 1 1,clip]{ex9-2-propagacao-x.png}
	\end{center}
\end{figure}

Alterando o ângulo $\theta$ da ponta de prova em relação ao centro do grid, 
podemos verificar como a reflexão da PML varia em relação ao ângulo de incidência. 
Variando as posições $p_x$ e $p_y$ da ponta de prova, tendo em mente que 
o ponto $\left( N_x / 2 \, , \, N_y / 2 \right)$ corresponde ao centro do grid, 
obtemos os seguintes resultados para o coeficiente de reflexão $R$,
exibidos na Tabela \ref{tab:reflexao-angulo}:

\begin{table}[h!]
	\caption{Variação do coeficiente de reflexão com o ângulo de incidência.}
	\vspace{0.5cm}
	\centering
		\begin{tabular}{c|c|c|c}

		$p_x$ & $p_y$ & $\theta$ & $R$ \\ \hline
		$N_x / 2 + 0$    & $N_y / 2 + 50$       & $0 ^\circ$       &    0.013      \\ \hline
		$N_x / 2 + 50$    & $N_y / 2 + 30$       & $30 ^\circ$       &    0.016      \\ \hline
		$N_x / 2 + 50$    & $N_y / 2 + 50$       & $45 ^\circ$       &    0.019      \\ \hline
		$N_x / 2 + 30$    & $N_y / 2 + 50$       & $60 ^\circ$       &    0.016      \\ \hline
		$N_x / 2 + 50$    & $N_y / 2 + 0$       & $90 ^\circ$       &    0.013      \\
		\end{tabular}
	\label{tab:reflexao-angulo}
\end{table}

Vemos que a reflexão é máxima quando a onda incidente 
está inclinada em 45 graus em relação à normal da borda, 
e é mínima quando a onda incide perpendicularmente 
(0 ou 90 graus). Assim, temos um desempenho pior da PML nas quinas do domínio, 
conforme esperado.

\chapter{Exercício 9.3}

Todo o código usado nesse exercício encontra-se no Anexo \ref{cap:ex9-3-codigo}.

Usamos os seguintes parâmetros de discretização:

\begin{itemize}
	\item $N_x = 120 \quad , \quad N_y = 60$
	\item $\Delta y = \Delta x = 5 \un{cm}$
	\item $\Delta t = \frac{0.9}{c \sqrt{(\Delta x)^{2} + (\Delta y)^{2}}} = 0.1 \un{ns} $ (condição CFL 2D)
	\item $N_t = 800$
\end{itemize}

A fonte aplicada no centro do domínio é uma função senoidal dada por 

\[ f(n) = \sin \left( 2\pi f_0 n \Delta t \right) \]

\noindent com $f_0 = 600 \un{MHz}$.

Os parâmetros da CPML adotados foram os seguintes:

\begin{itemize}
	\item $N_{pml} = 30 \quad , \quad m = 4$
	\item $R_{err} = 10^{-6}$
	\item $\eta = \sqrt{\frac{\mu}{\epsilon}}$
	\item $\sigma_{max} = -\frac{(m+1)\ln(R_{\mathrm{err}})}{2\eta N_{pml}} = 2.44$
\end{itemize}

A propagação obtida está exibida na Figura \ref{fig:ex9-3-propagacao}, 
onde podemos observar que o campo se propaga para fora do domínio sem reflexões significativas, 
indicando que os PMLs estão funcionando corretamente. 

\begin{figure}[htbp]
	\caption{\label{fig:ex9-3-propagacao}Propagação de $E_z$ ao longo do grid, sem reflexão nas bordas.}
	\begin{center}
        \includegraphics[width=\textwidth,trim=1 1 1 1,clip]{ex9-3-propagacao.png}
	\end{center}
\end{figure}

Olhando em um corte horizontal no eixo x passando pelo centro da cavidade,
podemos observar o decaimento do campo $E_z$ ao entrar nas regiões de PML, como mostra a Figura \ref{fig:ex9-3-propagacao-x}.

\begin{figure}[htbp]
	\caption{\label{fig:ex9-3-propagacao-x}Decaimento do campo $E_z$ ao longo do eixo $x$ ao entrar nas regiões de PML.}
	\begin{center}
        \includegraphics[width=0.8\textwidth,trim=1 1 1 1,clip]{ex9-3-propagacao-x.png}
	\end{center}
\end{figure}

\section{Item (a)}

Para calcular o coeficiente de reflexão, utilizamos o seguinte procedimento:

\begin{enumerate}
	\item Realiza a simulação completa com PML e coleta os pontos em uma ponta de prova perto do centro. 
	Esse é o campo total $E_{total}$
	\item Repete a simulação com as bordas muito longe (sem reflexão) e coleta os pontos em uma ponta de prova perto do centro. 
	Esse é o campo incidente $E_{inc}$
	\item Calcula o campo refletido como $E_{ref} = E_{total} - E_{inc}$
	\item Calcula o coeficiente de reflexão como $R = \frac{|E_{ref}|}{|E_{inc}|}$, usando decibeis (dB)
\end{enumerate}

Variando o valor do parâmetro $m$ da PML, obtemos os valores exibidos na Tabela \ref{tab:reflexao-m}.
Vericamos que o coeficiente de reflexão diminui à medida que aumentamos o valor de $m$, indicando uma melhor absorção da PML.

\begin{table}[h!]
	\caption{Variação do coeficiente de reflexão em função de $m$.}
	\vspace{0.5cm}
	\centering
		\begin{tabular}{c|c}

		$m$ & $R$ (dB)\\ \hline
		$3$    & $-22.21 $     \\ \hline
		$3.25$    & $-24.53$     \\ \hline
		$3.5$    & $-26.33$     \\ \hline
		$3.75$    & $-28.39$     \\ \hline
		$4$    & $-30.61$      \\
		\end{tabular}
	\label{tab:reflexao-m}
\end{table}

\section{Item (b)}

Implementando o grading geométrico para o perfil de $\sigma (x)$, obtemos 
o perfil de decaimento exibido ao longo do eixo $x$ exibido na Figura \ref{fig:ex9-3-propagacao-x-geometric}
obtemos os seguintes valores para o coeficiente de reflexão,
exibidos na Tabela \ref{tab:reflexao-geometrico}:

\begin{figure}[htbp]
	\caption{\label{fig:ex9-3-propagacao-x-geometric}Decaimento do campo $E_z$ ao longo do eixo $x$ ao entrar nas regiões de PML, 
	com função $\sigma$ com grading geométrico.}
	\begin{center}
        \includegraphics[width=0.8\textwidth,trim=1 1 1 1,clip]{ex9-3-propagacao-x-geometric.png}
	\end{center}
\end{figure}


\begin{table}[h!]
	\caption{Variação do coeficiente de reflexão em função de $g$.}
	\vspace{0.5cm}
	\centering
		\begin{tabular}{c|c}

		$g$ & $R$ (dB)\\ \hline
		$2.5$    & $-51.28 $     \\ \hline
		$3.0$    & $-52.71$     \\ \hline
		$3.5$    & $-52.98$      \\
		\end{tabular}
	\label{tab:reflexao-geometrico}
\end{table}

Verificamos que o coeficiente de reflexão é significativamente menor usando o grading geométrico 
do que quando usamos o grading polinomial, indicando uma melhor absorção da PML.
Além disso, o coeficiente se reduz ligeiramente ao aumentar o parâmetro $g$, mas a diferença não é tão grande quanto a
 variação do parâmetro $m$ no grading polinomial.


\chapter{Exercício 9.4}

Todo o código usado nesse exercício encontra-se no Anexo \ref{cap:ex9-4-codigo}.

Usamos os seguintes parâmetros de discretização (mesmos do exercício 9.3):

\begin{itemize}
	\item $N_x = 120 \quad , \quad N_y = 60$
	\item $\Delta y = \Delta x = 5 \un{cm}$
	\item $\Delta t = \frac{0.9}{c \sqrt{(\Delta x)^{2} + (\Delta y)^{2}}} = 0.1 \un{ns} $ (condição CFL 2D)
	\item $N_t = 800$
\end{itemize}

A fonte aplicada no centro do domínio é uma função senoidal dada por 

\[ f(n) = \sin \left( 2\pi f_0 n \Delta t \right) \]

\noindent com $f_0 = 600 \un{MHz}$.

\section{Item (a)}

Os parâmetros da UPML adotados foram os seguintes:

\begin{itemize}
	\item $N_{pml} = 30 \quad , \quad m = 3.75$
	\item $R_{err} = 10^{-8}$
	\item $\sigma_{max} = -\frac{(m+1)\ln(R_{\mathrm{err}})}{2 N_{pml} \cdot 377} = 0.116$
\end{itemize}

A propagação obtida está exibida na Figura \ref{fig:ex9-4-propagacao}, 
onde podemos observar que o campo se propaga para fora do domínio sem reflexões significativas, 
indicando que os PMLs estão funcionando corretamente. 

\begin{figure}[htbp]
	\caption{\label{fig:ex9-4-propagacao}Propagação de $E_z$ ao longo do grid, sem reflexão nas bordas.}
	\begin{center}
        \includegraphics[width=\textwidth,trim=1 1 1 1,clip]{ex9-4-propagacao.png}
	\end{center}
\end{figure}

Olhando em um corte horizontal no eixo x passando pelo centro da cavidade,
podemos observar o decaimento do campo $E_z$ ao entrar nas regiões de PML, como mostra a Figura \ref{fig:ex9-4-propagacao-x}.

\begin{figure}[htbp]
	\caption{\label{fig:ex9-4-propagacao-x}Decaimento do campo $E_z$ ao longo do eixo $x$ ao entrar nas regiões de PML.}
	\begin{center}
        \includegraphics[width=0.8\textwidth,trim=1 1 1 1,clip]{ex9-4-propagacao-x.png}
	\end{center}
\end{figure}

Usando o mesmo procedimento do exercício 9.3, obtemos um coeficiente de reflexão $R$ de
-33.67 dB, que é próximo do valor obtido no exercício anterior com a CPML. 

\section{Item (b)}

Com $f_0 = 100 \un{MHz}$ obtemos a propagação exibida na Figura \ref{fig:ex9-4-propagacao-100mhz}, 
e o novo coeficiente de reflexão aumenta consideralvemente para - 17.6 dB,
quase uma ordem de grandeza maior que $f_0 = 600 \un{MHz}$, indicando que a UPML
não possui um bom desempenho se estiver a poucos comprimentos de onda de distância da fonte.

\begin{figure}[htbp]
	\caption{\label{fig:ex9-4-propagacao-100mhz}Propagação de $E_z$ ao longo do grid para $f_0 = 100 \un{MHz}$.}
	\begin{center}
        \includegraphics[width=\textwidth,trim=1 1 1 1,clip]{ex9-4-propagacao-100mhz.png}
	\end{center}
\end{figure}

\chapter{Exercício 9.8}

Todo o código usado nesse exercício encontra-se no Anexo \ref{cap:ex9-8-codigo}.

Usamos os seguintes parâmetros de discretização (mesmos do exercício 9.3):

\begin{itemize}
	\item $N_x = 200 \quad , \quad N_y = 200$
	\item $\Delta y = \Delta x = 1 \un{mm}$
	\item $\Delta t = \frac{0.9}{c \sqrt{(\Delta x)^{2} + (\Delta y)^{2}}} = 2.12 \un{ps} $ (condição CFL 2D)
	\item $N_t = 800$
\end{itemize}

A fonte aplicada no centro do domínio é uma função gaussiana dada por 

\[ f(t) = e^{-\frac{1}{2}\left(\frac{n - 40}{S}\right)^2} \]

\noindent onde $S$ é um parâmetro que controla a abertura do pulso, estando relacionada com a frequência do sinal.

As progações obtidas com a UPML e CPML são semelhantes com os resultados
obtidos nos exercícios anteriores. O objetivo nesse exercício é variar o parâmetro $S$ e a localização da ponta 
de prova para verificar o desempenho da UPML e da CPML em função da frequência e do ângulo de incidência.

Como o pulso é gaussiano, para calcular o coeficiente de reflexão basta dividir a amplitude da onda incidente pela
onda que foi refletida, como visto na Figura \ref{fig:ex9-2-reflexao}.

A Tabela \ref{tab:resultado-9-8} sumariza o resultado obtido nesse exercício.

\begin{table}[h!]
	\caption{Desempenho da UPML e CPML em função da frequência do sinal da fonte e 
	do ângulo de incidência.}
	\vspace{0.5cm}
	\centering
		\begin{tabular}{c|c|c|c}

		$S$ & $\theta$ & $R_{CPML}$ (dB) & $R_{UPML}$ (dB) \\ \hline
		$3 \cdot \Delta t$    & $0$ & -54.23 & -41.64  \\ \hline
		$3 \cdot \Delta t$   & $15$ & -53.38 & -41.77     \\ \hline
		$3 \cdot \Delta t$    & $30$ & -53.4 & -42.65     \\ \hline
		$3 \cdot \Delta t$    & $45$ & -50.74 & -43.83     \\ \hline 
 		 \hline 
		$5 \cdot \Delta t$    & $0$ & -48.14 & -35.98    \\ \hline
		$5 \cdot \Delta t$   & $15$ & -47.33 & -35.38   \\ \hline
		$5 \cdot \Delta t$    & $30$ & -46.04 & -36.09    \\ \hline
		$5 \cdot \Delta t$    & $45$ & -44.3 & -37.23    \\ \hline
		 \hline 
		$7 \cdot \Delta t$    & $0$ & -43.93 & -31.35    \\ \hline
		$7 \cdot \Delta t$   & $15$ & -43.10 & -31.38    \\ \hline
		$7 \cdot \Delta t$    & $30$ & -41.87 & -31.86    \\ \hline
		$7 \cdot \Delta t$    & $45$ & -40.07 & -33.06   \\ \hline
		\end{tabular}
	\label{tab:resultado-9-8}
\end{table}

Em relação a frequência, vemos que tanto a CPML quanto a UPML apresentam 
melhor desempenho em altas frequências do que em baixas. Em todos os casos, 
contudo, o desempenho da CPML foi melhor do que o da UPML, apresentando um
coeficiente de reflexão cerca de 10 dB menor. 

Em relação ao ângulo de incidência, vemos que a CPML apresenta melhor desempenho 
quando $\theta = 0$, e conforme o ângulo aumenta, a CPML apresenta maior reflexão do sinal. 
Contudo, a UPML apresenta o comportamento inverso, sendo que era esperado que ela atenuasse melhor
os sinais incidentes em $\theta = 0$. De qualquer forma, a CPML também apresenta uma melhor reflexão
que a UPML para todos os ângulos de incidência. 

